% GENERATED FILE. Edit canonical lessons, then run npm run generate:books.
% canonical-source-hash: fnv1a64:0da26c2753c99b98
% canonical-lessons: FA-C483-sadeq, FA-C483-sedaqat, FA-C483-rok, FA-C483-dorugh-goftan, FA-C483-dorughgu

\chapter{Honesty and Lies}
\label{ch:fa-honesty-and-lies}
\hlchaptermodality{483}

\begin{chapteropening}
\textbf{By the end of this chapter:} \emph{I can talk about honesty and lying.}

Complete the last lesson of chapter 483: I can talk about honesty and lying.
\end{chapteropening}

\section[\texorpdfstring{\fa{صادق} (sâdeq)}{صادق (sâdeq)}]{\fa{صادق} (sâdeq) — honest, sincere}

\label{lesson:FA-C483-sadeq}



\begin{quote}
Before the new one: say the Persian for a grudge, then the Persian for ashamed.
\end{quote}

\subsection*{You'll want to know: \fa{صادق}}
\textbf{\fa{صادق}} (\emph{sâdeq}) — \textquotedblleft{}honest, sincere\textquotedblright{}.

\begin{quote}
\fa{دوست} \fa{من} \fa{همیشه} \fa{با} \fa{من} \fa{صادق} \fa{است}.
\end{quote}

\begin{quote}
\emph{dust-e man hamishe bâ man sâdeq ast.}
\end{quote}

\begin{quote}
My friend is always honest with me.
\end{quote}

\emph{bâ kasi sâdeq budan} — to be honest with someone.

\begin{cousinweb}
\emph{sâdeq} was borrowed from Arabic \emph{ṣādiq}, 'true, sincere; reliable' (EtymDB).
\end{cousinweb}

\begin{grammarlens}[title={What you've built}]
One more word for telling, explaining and describing at B1.
\end{grammarlens}

\subsection*{Your turn}
\begin{itemize}
\raggedright
  \item \emph{Say it:} \emph{sâdeq}
  \item \emph{Say it:} \emph{sâdeq}, once more
  \item \emph{Say it:} a sentence of your own with \emph{sâdeq}
  \item \emph{Recall:} say the Persian for a grudge, then the Persian for ashamed, then say \emph{sâdeq} again
\end{itemize}

\begin{tcolorbox}[breakable,colback=teal!4,colframe=teal!35!black,title={Before you move on}]
What does \fa{صادق} mean? (\textquotedblleft{}Honest, sincere\textquotedblright{}.) Say it once more.
\end{tcolorbox}
\section[\texorpdfstring{\fa{صداقت} (sedâqat)}{صداقت (sedâqat)}]{\fa{صداقت} (sedâqat) — honesty, sincerity}

\label{lesson:FA-C483-sedaqat}



\begin{quote}
Before the new one: say the Persian for ashamed, then the Persian for honest.
\end{quote}

\subsection*{You'll want to know: \fa{صداقت}}
\textbf{\fa{صداقت}} (\emph{sedâqat}) — \textquotedblleft{}honesty, sincerity\textquotedblright{}.

\begin{quote}
\fa{صداقت} \fa{برای} \fa{من} \fa{از} \fa{پول} \fa{مهمتر} \fa{است}.
\end{quote}

\begin{quote}
\emph{sedâqat barâ-ye man az pul mohemtar ast.}
\end{quote}

\begin{quote}
Honesty matters more to me than money.
\end{quote}

A noun; it shares its Arabic root with \emph{sâdeq} 'honest'.

\begin{cousinweb}
\emph{sedâqat} is a loan, usually traced to Arabic \emph{ṣadāqa}, a word from the same root as \emph{ṣādiq}, 'true, sincere'. In Arabic \emph{ṣadāqa} means 'friendship'; Persian took it as 'honesty'. EtymDB has no entry for it.
\end{cousinweb}

\begin{grammarlens}[title={What you've built}]
One more word for telling, explaining and describing at B1.
\end{grammarlens}

\subsection*{Your turn}
\begin{itemize}
\raggedright
  \item \emph{Say it:} \emph{sedâqat}
  \item \emph{Say it:} \emph{sedâqat}, once more
  \item \emph{Say it:} a sentence of your own with \emph{sedâqat}
  \item \emph{Recall:} say the Persian for ashamed, then the Persian for honest, then say \emph{sedâqat} again
\end{itemize}

\begin{tcolorbox}[breakable,colback=teal!4,colframe=teal!35!black,title={Before you move on}]
What does \fa{صداقت} mean? (\textquotedblleft{}Honesty, sincerity\textquotedblright{}.) Say it once more.
\end{tcolorbox}
\section[\texorpdfstring{\fa{رک} (rok)}{رک (rok)}]{\fa{رک} (rok) — frank, blunt}

\label{lesson:FA-C483-rok}



\begin{quote}
Before the new one: say the Persian for honest, then the Persian for honesty.
\end{quote}

\subsection*{You'll want to know: \fa{رک}}
\textbf{\fa{رک}} (\emph{rok}) — \textquotedblleft{}frank, blunt\textquotedblright{}.

\begin{quote}
\fa{او} \fa{خیلی} \fa{رک} \fa{است} \fa{و} \fa{نظرش} \fa{را} \fa{میگوید}.
\end{quote}

\begin{quote}
\emph{u kheyli rok ast va nazarash râ miguyad.}
\end{quote}

\begin{quote}
She's very frank and says what she thinks.
\end{quote}

Common in speech; also \emph{rok-o-râst} 'straight out'.

\begin{cousinweb}
\emph{rok} is an everyday Persian word for frank, but no source used here gives its history, and its origin is not settled.
\end{cousinweb}

\begin{grammarlens}[title={What you've built}]
One more word for telling, explaining and describing at B1.
\end{grammarlens}

\subsection*{Your turn}
\begin{itemize}
\raggedright
  \item \emph{Say it:} \emph{rok}
  \item \emph{Say it:} \emph{rok}, once more
  \item \emph{Say it:} a sentence of your own with \emph{rok}
  \item \emph{Recall:} say the Persian for honest, then the Persian for honesty, then say \emph{rok} again
\end{itemize}

\begin{tcolorbox}[breakable,colback=teal!4,colframe=teal!35!black,title={Before you move on}]
What does \fa{رک} mean? (\textquotedblleft{}Frank, blunt\textquotedblright{}.) Say it once more.
\end{tcolorbox}
\section[\texorpdfstring{\fa{دروغ} \fa{گفتن} (dorugh goftan)}{دروغ گفتن (dorugh goftan)}]{\fa{دروغ} \fa{گفتن} (dorugh goftan) — to lie, to tell a lie}

\label{lesson:FA-C483-dorugh-goftan}



\begin{quote}
Before the new one: say the Persian for honesty, then the Persian for frank.
\end{quote}

\subsection*{You'll want to know: \fa{دروغ} \fa{گفتن}}
\textbf{\fa{دروغ} \fa{گفتن}} (\emph{dorugh goftan}) — \textquotedblleft{}to lie, to tell a lie\textquotedblright{}.

\begin{quote}
\fa{بچه} \fa{ها} \fa{گاهی} \fa{به} \fa{پدر} \fa{و} \fa{مادرشان} \fa{دروغ} \fa{میگویند}.
\end{quote}

\begin{quote}
\emph{bachchehâ gâhi be pedar-o mâdarashân dorugh miguyand.}
\end{quote}

\begin{quote}
Children sometimes lie to their parents.
\end{quote}

\emph{dorugh} 'a lie' + \emph{goftan} 'to say'; present stem \emph{gu-}. The person lied to takes \emph{be}. The opposite is \emph{râst goftan}.

\begin{cousinweb}
\emph{dorugh} is inherited: EtymDB traces it to Old Persian \emph{drauga}, 'lie', and on to a Proto-Indo-European root for deceiving, and lists German \emph{Betrug}, 'deceit', as a cousin (EtymDB). \emph{goftan}, 'to say', is inherited too.
\end{cousinweb}

\begin{grammarlens}[title={What you've built}]
One more word for telling, explaining and describing at B1.
\end{grammarlens}

\subsection*{Your turn}
\begin{itemize}
\raggedright
  \item \emph{Say it:} \emph{dorugh goftan}
  \item \emph{Say it:} \emph{dorugh goftan}, once more
  \item \emph{Say it:} a sentence of your own with \emph{dorugh goftan}
  \item \emph{Recall:} say the Persian for honesty, then the Persian for frank, then say \emph{dorugh goftan} again
\end{itemize}

\begin{tcolorbox}[breakable,colback=teal!4,colframe=teal!35!black,title={Before you move on}]
What does \fa{دروغ} \fa{گفتن} mean? (\textquotedblleft{}To lie, to tell a lie\textquotedblright{}.) Say it once more.
\end{tcolorbox}
\section[\texorpdfstring{\fa{دروغگو} (dorughgu)}{دروغگو (dorughgu)}]{\fa{دروغگو} (dorughgu) — a liar}

\label{lesson:FA-C483-dorughgu}



\begin{quote}
Before the new one: say the Persian for frank, then the Persian for to lie.
\end{quote}

\subsection*{You'll want to know: \fa{دروغگو}}
\textbf{\fa{دروغگو}} (\emph{dorughgu}) — \textquotedblleft{}a liar\textquotedblright{}.

\begin{quote}
\fa{هیچکس} \fa{حرف} \fa{یک} \fa{دروغگو} \fa{را} \fa{باور} \fa{نمیکند}.
\end{quote}

\begin{quote}
\emph{hichkas harf-e yek dorughgu râ bâvar nemikonad.}
\end{quote}

\begin{quote}
Nobody believes a liar.
\end{quote}

\emph{dorugh} 'lie' + \emph{gu} 'saying' (stem of \emph{goftan}). Plural \emph{dorughguhâ}.

\begin{cousinweb}
\emph{dorughgu} is built on \emph{dorugh}, 'lie', which EtymDB traces to Old Persian \emph{drauga} (EtymDB), and \emph{gu}, the present stem of \emph{goftan}, 'to say': a lie-teller.
\end{cousinweb}

\begin{grammarlens}[title={What you've built}]
One more word for telling, explaining and describing at B1.
\end{grammarlens}

\subsection*{Your turn}
\begin{itemize}
\raggedright
  \item \emph{Say it:} \emph{dorughgu}
  \item \emph{Say it:} \emph{dorughgu}, once more
  \item \emph{Say it:} a sentence of your own with \emph{dorughgu}
  \item \emph{Recall:} say the Persian for frank, then the Persian for to lie, then say \emph{dorughgu} again
\end{itemize}

\begin{tcolorbox}[breakable,colback=teal!4,colframe=teal!35!black,title={Before you move on}]
What does \fa{دروغگو} mean? (\textquotedblleft{}A liar\textquotedblright{}.) Say it once more.
\end{tcolorbox}
